\documentclass[10pt]{extarticle}
\usepackage[a4paper,margin=1.5cm]{geometry}
\usepackage[T1]{fontenc}
\usepackage{amsmath,amssymb}
\usepackage{graphicx}
\usepackage[export]{adjustbox}
\usepackage{multicol}
\usepackage{enumitem}
\usepackage{booktabs,array,tabularx}
\setlength{\columnsep}{0.7cm}
\setlength{\parindent}{0pt}
\setlength{\parskip}{3pt}
\pagestyle{empty}
\makeatletter
% figure registry (figdefs.tex is generated by the extract step)
\newcommand{\qfigdef}[4]{\@namedef{qf@file@#1}{#2}\@namedef{qf@w@#1}{#3}\@namedef{qf@wide@#1}{#4}}
\newcommand{\qf@check}[1]{\@ifundefined{qf@file@#1}{\PackageError{score2}{Unknown figure id #1}{}}{}}
\newcommand{\qf@img}[4][]{\includegraphics[width=#2\dimexpr\@nameuse{qf@w@#4}pt\relax,max width=#3,max height=0.45\textheight,#1]{\@nameuse{qf@file@#4}}}
% Full-width questions: finalize wraps every question that contains a wide figure in \qfullbegin...\qfullend,
% so the WHOLE question (stem, figures, options) is set across the page in reading order, then the two columns
% resume. The question's head - its number, stem and first figure (\qfullsplit marks the end of the head) -
% is kept on one page, so a number is never left at a page foot with its figure overleaf; the rest may break
% across pages normally, so no big gaps are left. Without a \qfullsplit, the whole question is kept together
% if it is under 45% of a page. Never write these yourself.
\newif\ifqf@full
\newif\ifqf@boxed
\newsavebox{\qf@box}
\newcommand{\qf@open}{\global\qf@fulltrue\par\global\qf@boxedtrue
  \setbox\qf@box\vbox\bgroup\hsize\textwidth\linewidth\textwidth}
\newcommand{\qf@emit}[1]{\par\egroup\global\qf@boxedfalse
  \ifdim\dimexpr\ht\qf@box+\dp\qf@box\relax<#1\textheight\noindent\usebox{\qf@box}\par\else\unvbox\qf@box\fi}
\newcommand{\qfullsplit}{\ifqf@boxed\qf@emit{0.85}\fi}
\newcommand{\qf@close}{\ifqf@boxed\qf@emit{0.45}\fi\global\qf@fullfalse}
\newcommand{\qfullbegin}{\end{multicols}\qf@open}
\newcommand{\qfullend}{\qf@close\begin{multicols}{2}}
\newcommand{\qfullstart}{\qf@open}
\newcommand{\qfulldone}{\qf@close}
% \qfig[scale]{id}: block figure, centred, directly where it is written (size set by the tool).
% Fallback if a wide figure is not inside a full-width question: break out of the columns in place - no floats.
\newcommand{\qfig}[2][1]{\qf@check{#2}%
  \ifqf@full\qf@blk{#1}{#2}\else
  \ifnum\@nameuse{qf@wide@#2}=1 \expandafter\@firstoftwo\else\expandafter\@secondoftwo\fi
  {\end{multicols}\begin{center}\qf@img{#1}{\textwidth}{#2}\end{center}\begin{multicols}{2}}%
  {\qf@blk{#1}{#2}}\fi}
\newcommand{\qf@blk}[2]{\par\smallskip{\centering\qf@img{#1}{\linewidth}{#2}\par}\smallskip}
% \qfigi[scale]{id}: inline figure (inside option lists / tables / minipages); never breaks columns.
\newcommand{\qfigi}[2][1]{\qf@check{#2}\qf@img[valign=c]{#1}{\linewidth}{#2}}
\makeatother
% structure macros - lint reads these to check completeness and figure placement
\newcommand{\papersec}[1]{\par\bigskip{\centering\Large\bfseries PAPER-#1\par}\medskip}
\newcommand{\partsec}[2]{\par\medskip\noindent{\large\bfseries PART-#1: #2}\par\smallskip}
\newcommand{\sect}[1]{\par\smallskip\noindent\textbf{#1}\par}
\newcommand{\qn}[1]{\par\smallskip\noindent\textbf{#1.}~}
\newcommand{\paraq}[2]{\par\smallskip\noindent\textbf{Paragraph for Questions #1 and #2}\par}
\newenvironment{qoptions}{\begin{itemize}[leftmargin=1.8em,itemsep=1pt,topsep=2pt,label={}]}{\end{itemize}}
\newcommand{\opt}[2]{\item[(#1)] #2}
% figure registry (generated by score2_tools.py finalize)
\qfigdef{p1_phy_s0_q02_1}{SRGBT - Test 01 (25 Jul 2021)-img/p1_phy_s0_q02_1.png}{102.2}{0}
\qfigdef{p1_phy_s0_q04_1}{SRGBT - Test 01 (25 Jul 2021)-img/p1_phy_s0_q04_1.png}{100.3}{0}
\qfigdef{p1_phy_s0_q05_1}{SRGBT - Test 01 (25 Jul 2021)-img/p1_phy_s0_q05_1.png}{63.7}{0}
\qfigdef{p1_phy_s0_q06_1}{SRGBT - Test 01 (25 Jul 2021)-img/p1_phy_s0_q06_1.png}{92.0}{0}
\qfigdef{p1_phy_s0_q08_1}{SRGBT - Test 01 (25 Jul 2021)-img/p1_phy_s0_q08_1.png}{96.8}{0}
\qfigdef{p1_phy_s0_q09_1}{SRGBT - Test 01 (25 Jul 2021)-img/p1_phy_s0_q09_1.png}{62.0}{0}
\qfigdef{p1_phy_s0_q10_1}{SRGBT - Test 01 (25 Jul 2021)-img/p1_phy_s0_q10_1.png}{100.3}{0}
\qfigdef{p1_phy_s0_q13_1}{SRGBT - Test 01 (25 Jul 2021)-img/p1_phy_s0_q13_1.png}{71.2}{0}
\qfigdef{p1_phy_s0_q15_2}{SRGBT - Test 01 (25 Jul 2021)-img/p1_phy_s0_q15_2.png}{124.3}{0}
\qfigdef{p1_phy_s0_q17_1}{SRGBT - Test 01 (25 Jul 2021)-img/p1_phy_s0_q17_1.png}{82.9}{0}
\raggedcolumns
\makeatletter\@beginparpenalty=10000\makeatother
\newcolumntype{L}[1]{>{\raggedright\arraybackslash}p{#1}}
\begin{document}
\begin{center}{\Large\bfseries SRGBT --- Test 01 (25 Jul 2021)}\end{center}
\begin{multicols}{2}
\raggedcolumns
\papersec{1}
\partsec{1}{PHYSICS}
\qn{1} In a football game, a player wants to throw a ball to his one of teammate, who is running on the field. Taken thrower position as origin and receiver initial position at $2\hat i+3\hat j$, where $\hat i$ \& $\hat j$ are in the plane of field. In subsequent run receiver displacement is $5\hat i$ then $8\hat j$, then $-2\hat i+4\hat j$ then $-6\hat j$. How far is receiver from thrower? (all displacements are in meter)
\begin{qoptions}
\opt{A}{$\sqrt{106}$}
\opt{B}{$\sqrt{110}$}
\opt{C}{$\sqrt{103}$}
\opt{D}{$\sqrt{112}$}
\end{qoptions}
\qn{2} A particle is projected with speed u from origin at an angle 45$^\circ$ with positive x-axis. There are two parallel imaginary lines L$_1$ and L$_2$ in the x-y plane as shown in figure. Line L$_1$ passes through $\left(\dfrac R4,0\right)$ and line L$_2$ passes through $\left(\dfrac R2,0\right)$. Lines are inclined at 45$^\circ$ with x-axis. Acceleration due to gravity is `g' along negative y-axis. [where R = $\dfrac{\mathrm u^2}{\mathrm g}$].\par
\qfig{p1_phy_s0_q02_1}
The angle $\theta$ between acceleration and velocity of particle between two instant of crossing the lines L$_1$ and L$_2$ is always :
\begin{qoptions}
\opt{A}{$\theta>45^\circ$}
\opt{B}{$\theta<45^\circ$}
\opt{C}{$\theta<90^\circ$}
\opt{D}{$\theta>90^\circ$}
\end{qoptions}
\qn{3} The current velocity of a river grows in proportion to the distance from its bank and reaches the maximum value v$_0$ in the middle. Near the banks the velocity is zero. A boat is moving along the river in such a manner that it is always perpendicular to the current and the speed of the boat in still water is u. Find the distance through which the boat crossing the river will be carried away by the current if the width of the river is c.
\begin{qoptions}
\opt{A}{$\dfrac{\mathrm{cv}_0}{\mathrm u}$}
\opt{B}{$\dfrac{\mathrm{cv}_0}{2\mathrm u}$}
\opt{C}{$\dfrac{\mathrm{cv}_0}{4\mathrm u}$}
\opt{D}{$\dfrac{\mathrm{cv}_0}{8\mathrm u}$}
\end{qoptions}
\qn{4} A bus is moving on a highway with a speed of v$_1$ = 20 m/s. A man is at a perpendicular distance of a = 60 m from the highway and b = 400 m from the bus. In order to catch the bus suppose he runs at an angle $\alpha$ from the line of joining him to the bus in initial condition. The man can run at a speed of v$_2$ = 6 m/s. Then the possible values of $\alpha$ so that he can reach any point on the highway at the same time as the bus or before it. Highway at the same time as the bus or before it.\par
\qfig{p1_phy_s0_q04_1}
\begin{qoptions}
\opt{A}{160$^\circ$}
\opt{B}{180$^\circ$}
\opt{C}{15$^\circ$}
\opt{D}{120$^\circ$}
\end{qoptions}
\qn{5} A particle starts moving with initial velocity 3 m/s along x-axis from origin. Its acceleration is varying with x in parabolic nature as shown in figure. At x = $\sqrt3$ m tangent to the graph makes an angle 60$^\circ$ with positive x-axis as shown in diagram. Then at x = $\sqrt3$\par
\qfig{p1_phy_s0_q05_1}
\begin{qoptions}
\opt{A}{v = $\sqrt{(\sqrt3+9)}$ m / s}
\opt{B}{a = 1.5 m/s$^2$}
\opt{C}{v = $\sqrt{12}$ m/s}
\opt{D}{a = 3 m/s$^2$}
\end{qoptions}
\qn{6} A small ball is moving with the speed of 6 m/s on horizontal surface falls in a vertical trench 4m-wide and 20m deep. Consider all collisions elastic and friction is absent every where. choose the correct statement(s).\par
\qfig{p1_phy_s0_q06_1}
\begin{qoptions}
\opt{A}{The ball will come out of the trench at A}
\opt{B}{The ball will come out of the trench at B}
\opt{C}{The ball takes time 4s time and comes out}
\opt{D}{The ball never comes out of it}
\end{qoptions}
\qn{7} A cylindrical pipe of radius r is at rest. A frog sitting on the horizontal ground jumps off and passes over the pipe touching it only at the top. Denoting air-time of the frog by T, horiozntal range of the jump by R and acceleration due to gravity by g, which of the following conclusions can you make ?
\begin{qoptions}
\opt{A}{T = 4$\sqrt{\dfrac rg}$}
\opt{B}{T = 2$\sqrt{\dfrac rg}$}
\opt{C}{R = 2$\left(\sqrt{\mathrm{gr}}\right)\sqrt{\dfrac rg}$}
\opt{D}{R = 4$\left(\sqrt{\mathrm{gr}}\right)\sqrt{\dfrac rg}$}
\end{qoptions}
\qn{8} A smooth inclined plane has an angle of inclination $\theta$ from the horizontal. From point A on the inclined plane, body 1 is projected horizontally with speed u and from point B on the inclined plane, body 2 is released from rest. Distance AB is equal to $\ell$. If the magnitude of relative velocity of body 2 w.r.t. body 1 and magnitude of relative acceleration of body 2 w.r.t. body 1 at time t are v and `a', respectively, then which of the following is/are true?\par
\qfig{p1_phy_s0_q08_1}
\begin{qoptions}
\opt{A}{v = $\sqrt{\mathrm u^2-2\mathrm{ug}\sin\theta\cos\theta\,\mathrm t+\mathrm g^2\cos^2\theta\,\mathrm t^2}$}
\opt{B}{a = zero}
\opt{C}{a = g cos $\theta$}
\opt{D}{v = $\sqrt{\left(\mathrm g\sin^2\theta\,\mathrm t-\mathrm u\right)^2+\left(\mathrm g\sin\theta\cos\theta\,\mathrm t-\mathrm{gt}\right)^2}$}
\end{qoptions}
\qn{9} A 600 m wide river flows towards South at 4 m/s. A small motor boat travels at 5 m/s in still water \& points in such a direction so that it will travel towards East relative to land.\par
\qfig{p1_phy_s0_q09_1}
\begin{qoptions}
\opt{A}{The speed of boat relative to land is 3 m/s.}
\opt{B}{The speed of boat relative to land is $\sqrt{41}$ m/s.}
\opt{C}{The time taken by boat to cross the river is nearly 200 second.}
\opt{D}{The time taken by boat to cross the river is nearly 100 second.}
\end{qoptions}
\qn{10} A particle of mass m is projected up from the bottom of an smooth inclined plane with initial velocity v$_0$ at angle 45$^\circ$ with an inclined plane of inclination 30$^\circ$ as shown in figure. At the same time a small block of same mass m is released from rest at a height h, the particle hits the block at some point on inclined plane. (take g = 10 m/s$^2$) (size of the block is negligible)\par
\qfig{p1_phy_s0_q10_1}
\begin{qoptions}
\opt{A}{The value of height h is $\dfrac{\mathrm v_0^2}{\sqrt3\mathrm g}$.}
\opt{B}{The value of height h is $\dfrac{\sqrt3\mathrm v_0^2}{\mathrm g}$.}
\opt{C}{If the particle stick to the block after collision, the velocity of block parallel to the inclined plane just after collision will be $\dfrac{\mathrm v_0}{\sqrt2}$.}
\opt{D}{If the particle stick to the block after collision, the velocity of block parallel to the inclined plane just after collision will be $\dfrac{\mathrm v_0}{2\sqrt2}\left(\dfrac4{\sqrt3}-1\right)$.}
\end{qoptions}
\qn{11} A child throws a ball so as to clear a wall of height h and at a distance x from it. Find the minimum speed required for clearing the wall is $\sqrt{\mathrm g\left(\mathrm{ah}+\sqrt{\mathrm{bh}^2+\mathrm{cx}^2}\right)}$ then find a + b + c ?
\qn{12} Two particles are projected from a point at the same instant with velocities whose horizontal and vertical components are u$_1$, v$_1$ and u$_2$, v$_2$ respectively. Prove that the interval between their passing through the other common point of their path is $\dfrac{\mathrm a(\mathrm{bv}_1\mathrm u_2+\mathrm{cv}_2\mathrm u_1)}{\mathrm g(\mathrm u_1+\mathrm u_2)\mathrm d}$. Find a + b + c + d ?
\qn{13} A particle is released from the top of an equilateral triangular groove which is meet in the plane of an incline plane of inclination $\theta$ = 53$^\circ$. The corners of triangular groove are curved to neglect impulse. Side of triangle is 8$\sqrt3$ m. Calculate the time of a complete round trip (at corners there is no change in speed).\par
\qfig{p1_phy_s0_q13_1}
\qn{14} Two tourists who are at a distance of 40 km from their camp must reach it together in shortest possible time. They have one bicycle which they decide to use in turn. One of them started walking at a speed of 5 km/hr and the other rod of on the bicycle at a speed of 15 km/hr. The tourists agreed to leave bicycle at intermediate points between which one walks and other ride. How long bicycle remain unused (in minutes) :
\qn{15} Figure shows top view of a boat moving with velocity $\vec{\mathrm v}_\mathrm r=\left(8\hat i+6\hat j\right)$ m/s relative to water. Water flows with velocity $\vec{\mathrm v}_\mathrm w=2\hat i$ m/s relative to ground. A motor cyclist moves along a road running parallel to bank with velocity $\vec{\mathrm v}_\mathrm m=20\hat i$ m/s. An apple thrown from boat is caught by motor cyclist in same horizontal plane after 4 sec of projection. If apple is thrown with velocity $\mathrm v_\mathrm x\hat i+\mathrm v_\mathrm y\hat j+\mathrm v_\mathrm z\hat k$ relative to boat fill value of $\dfrac{|\vec{\mathrm v}_\mathrm x|+|\vec{\mathrm v}_\mathrm y|+|\vec{\mathrm v}_\mathrm z|}{13}$.\par
\qfig{p1_phy_s0_q15_2}
\qn{16} A boy crosses a river twice on a straight path at an angle $\phi$ = 60$^\circ$ with the downstream direction, first time in two minutes andsecond time in four minutes. If his speed relative to river current is v$_\mathrm{b/w}$ = $\sqrt3$ m/s in both the attempts, find speed of the river current.
\qn{17} Two point A and B are located on floor of a room as shown in figure. The room has three walls 1, 2 and 3. Perpendicular distance of point A from wall 1, 2, and 3 is 0.4 m, 2m and 3m respectively. While that of point B is 3.2 m, 1m and 0.2 m respectively. A man has to move from point A to B touching the walls 1, 2 and 3 at least once. Minimum distance (in meter) traveled by man will be :\par
\qfig{p1_phy_s0_q17_1}
\qn{18} The x-coordinate of a particle moving on x-axis is given by x = 3 sin 100t + 8 cos$^2$50t, where x is in cm and t is time in seconds. Find the maximum distance the particle can move from the origin (in cm).
\raggedcolumns
\partsec{2}{CHEMISTRY}
\qn{19} Calculate the sum of stoichiometirc coefficients of products in minimum integral balanced following chemical reaction.\par
a CrCl$_3$ + b H$_2$O$_2$ + c NaOH $\longrightarrow$ x Na$_2$CrO$_4$ + y NaCl + z H$_2$O
\begin{qoptions}
\opt{A}{8}
\opt{B}{16}
\opt{C}{45}
\opt{D}{32}
\end{qoptions}
\qn{20} 100 m$l$ of a gaseous mixture containing Ne, CO$_2$ \& H$_2$ on complete combustion in just sufficient amount of O$_2$ showed contraction of 60 ml at STP. When the resulting gases were passed through KOH solution, volume reduces by 40\%. The volume ratio of V$_{\mathrm{CO}_2}$ : V$_\mathrm{Ne}$ : V$_{\mathrm H_2}$ in original mixture is.
\begin{qoptions}
\opt{A}{10 : 9 : 2}
\opt{B}{6 : 9 : 10}
\opt{C}{7 : 3 : 10}
\opt{D}{None of these}
\end{qoptions}
\qn{21} The molality of pure gas A is $\dfrac{50}{2.4}$ and that of pure B is $\dfrac{100}{3.2}$. Calculate molarity of A in a gaseous solution of A and B at 8 atm pressure and $\dfrac{200}{0.821}$ K. if the vapour density of the gaseous mixture at 10 atm and 400 K is 18.
\begin{qoptions}
\opt{A}{0.4}
\opt{B}{0.1}
\opt{C}{0.3}
\opt{D}{0.2}
\end{qoptions}
\qn{22} 100 ml of a sample of hard water requires 25.1 ml of 0.02 N H$_2$SO$_4$ for complete reaction. The hardness of water (density = 1 gm/ml) is :
\begin{qoptions}
\opt{A}{200 ppm}
\opt{B}{251 ppm}
\opt{C}{261 ppm}
\opt{D}{288 ppm}
\end{qoptions}
\qn{23} 40 ml, 0.05 M solution of sodium sesquicarbonate dihydrate [Na$_2$CO$_3$.NaHCO$_3$.2H$_2$O] is titrated against 0.05 M HCl solution, x ml of acid is required to reach the phenolphthalein end point, while y ml. of same acid was required when methyl orange indicator was used in a separate titration. Which of the following statements are correct ?
\begin{qoptions}
\opt{A}{y $-$ x = 80 ml}
\opt{B}{y + x = 160 ml}
\opt{C}{If the titration is started with phenolphthalein indicator and methyl orange is added at the end point, 2x ml. HCl would be required further to reach end point.}
\opt{D}{y + x = 120 ml.}
\end{qoptions}
\qn{24} Which of the following options represent concentration terms which are independant of temperature ?
\begin{qoptions}
\opt{A}{\% w/w, mole fraction, molality}
\opt{B}{Volume strength of H$_2$O$_2$, normality, ppm}
\opt{C}{Molality, mass fraction, \% labelling}
\opt{D}{Molarity, \% w/v, ,mole fraction}
\end{qoptions}
\qn{25} A definite volume of pure ammonia (NH$_3$) gas is passed through a series of electric sparks by which the volume becomes 90 ml. The increase in volume is due to formation of nitrogen (N$_2$) and Hydrogen (H$_2$)gases. All the gases finally present are washed with dilute sulphuric acid solution, by which the volume of gases become 80 ml. All the volumes are measured at the same temperature and pressure. Which of the following statement(s) is/are \textbf{correct} regarding the original ammonia sample?
\begin{qoptions}
\opt{A}{The volume of NH$_3$ gas taken was 40 ml.}
\opt{B}{The volume of NH$_3$ gas taken was 50 ml.}
\opt{C}{Only 80\% of NH$_3$ gas decomposed into N$_2$ and H$_2$ gases.}
\opt{D}{Only 20\% of NH$_3$ gas decomposed into N$_2$ and H$_2$ gases.}
\end{qoptions}
\qn{26} Which of the following option/s is/are correct w.r.t equivalent mass (EM).
\begin{qoptions}
\opt{A}{Equivalent mass of HCl in its reaction with KMnO$_4$ to give MnCl$_2$ \& Cl$_2$ is equal to 36.5}
\opt{B}{Equivalent mass of MnO$_2$ when it converts to Mn$_2$O$_7$ \& Mn$^{2+}$ in a molar ratio of 2 : 1 is $\dfrac56$ times of its molecular mass}
\opt{C}{Equivalent mass of NaOH in its reaction with S$_8$ to give Na$_2$SO$_4$ \& Na$_2$S is $\dfrac43$ times its molecular mass.}
\opt{D}{Equivalent mass of K$_4$[Fe(CN)$_6$] when it reacts with FeCl$_3$ to form KFe[Fe(CN)$_6$] is $\dfrac13$ of its molecular mass.}
\end{qoptions}
\qn{27} For the reaction\par
CrO$_5$ + H$_2$SO$_4$ $\longrightarrow$ Cr$_2$(SO$_4$)$_3$ + H$_2$O + O$_2$\par
Incorrect statement is/are :
\begin{qoptions}
\opt{A}{Oxidation state of `Cr' in CrO$_5$ is +8.}
\opt{B}{Oxidation of one of element in CrO$_5$ takes place.}
\opt{C}{Reduction of one of element in CrO$_5$ takes place.}
\opt{D}{Oxidation state of `Cr' in CrO$_5$ is +10.}
\end{qoptions}
\qn{28} 1 mole of a mixture of CO \& CO$_2$ requires exactly 1 litre solution of 1 M NaOH for complete neutralisation. If CO present in mixture is now converted to CO$_2$ and again the mixture is treated with NaOH, then after this conversion :
\begin{qoptions}
\opt{A}{moles of CO$_2$ present initially in mixture = 1.}
\opt{B}{2 litre NaOH solution of 1 M is more required for neutralisation.}
\opt{C}{2 litre solution of $\dfrac{\mathrm M}2$ is more required for neutralisation.}
\opt{D}{56 gm KOH in aqueous solution is required more for neutralisation.}
\end{qoptions}
\qn{29} Given balanced chemical equation for oxidation of phosphorus (III) sulphide by nitric acid. The products include NO and SO$_2$. Find the value of a, b, c, d, e in the following equation.\par
a P$_4$S$_6$ + 44H$^+$ + b NO$_3^-$ $\to$ cNO + dH$_3$PO$_4$ + eSO$_2$ + 4H$_2$O\par
find the value of (a + b + c + d + e)
\qn{30} 50 ml of gaseous mixture of acetylene and ethylene if taken in a ratio of a:b requires 700 ml of air containing 20\% by volume O$_2$ for complete combustion. Calculate the volume of air required for complete combustion of a mixture having ratio b : a.
\qn{31} Composition of Oleum \hfill Solution Code\par
{\small\setlength{\tabcolsep}{3pt}\noindent\begin{tabular}{ll@{\quad}c}
50 gm SO$_3$ ; & 120 gm H$_2$SO$_4$ & 5 \\
230 gm SO$_3$ ; & 100 gm H$_2$SO$_4$ & 6 \\
80 gm SO$_3$ ; & 98 gm H$_2$SO$_4$ & 2 \\
133 gm SO$_3$ ; & 86 gm H$_2$SO$_4$ & 3 \\
93 gm SO$_3$ ; & 93 gm H$_2$SO$_4$ & 1 \\
400 gm SO$_3$ ; & 500 gm H$_2$SO$_4$ & 4 \\
\end{tabular}}\par
Mention the solution code of the Oleum having maximum \% labelling.
\qn{32} Consider the production of tetraethyl lead according to the following reaction.\par
4C$_2$H$_5$Cl + 4Na - Pb $\longrightarrow$ (C$_2$H$_5$)$_4$Pb + 4NaCl + 3Pb\par
How much kilograms of ethyl chloride is required to produce enough tetraethyl lead (density $\dfrac{12.96}{1.29}$ g / ml) needed for 500 litre of aviation fuel using 2 ml of tetraethyl lead per litre of fuel ? (Pb = 208)
\qn{33} 200 mL of 1 M HCl solution is mixed with 800 mL of 2M HCl solution. Calculate the molarity of final solution.
\qn{34} A mixture of formic acid and oxalic acid is heated in the presence of H$_2$SO$_4$.\par
HCOOH $\longrightarrow$ H$_2$O + CO\par
H$_2$C$_2$O$_4$ $\longrightarrow$ H$_2$O + CO$_2$ + CO\par
The produced gaseous mixture is collected and cooled. On its treatment with KOH solution, the volume of gaseous mixture decreases by one-tenth. Calculate the molar ratio of the two acid in original mixturel.[If the simplest ratio is a : b then find a + b].
\qn{35} In the following reactions if initially 100 moles of A were taken then calculate sum of moles of D and C after the reaction.\par
A $\xrightarrow{40\%}$ 3B + D\par
2B $\xrightarrow{50\%}$ 4C
\qn{36} If 246 g of a triatomic ideal gas having only one type of atoms is found to occupy 44.8 L at 2 atm and 546 K then what will be the atomic weight of the element ?
\raggedcolumns
\partsec{3}{MATHEMATICS}
\qn{37} If $\sin^2\mathrm x-2(\mathrm a+1)\sin\mathrm x+\left(\mathrm a^2+2\mathrm a-4\right)>0$ for all $\mathrm x\in\mathrm R$, then the complete set of values of a is $(-\infty,\mathrm m)\cup(\mathrm n,\infty)$.\par
Find the value of [n] $-$ [m]\par
( [.] denotes greatest integer function)
\begin{qoptions}
\opt{A}{5}
\opt{B}{6}
\opt{C}{7}
\opt{D}{8}
\end{qoptions}
\qn{38} Let f(x) be a polynomial function of degree at most 5 which leaves remainder $-$ 1 and 1 upon divison by (x $-$ 1)$^3$ and (x + 1)$^3$ respectively . The value of f(2) is equal to
\begin{qoptions}
\opt{A}{$\dfrac{46}{15}$}
\opt{B}{$\dfrac{-23}8$}
\opt{C}{$\dfrac{-23}4$}
\opt{D}{none of these}
\end{qoptions}
\qn{39} Sum of all the integral solutions of the equation\par
{\footnotesize$\sqrt{\left(\log_3\sqrt[3]{3\mathrm x}+\log_\mathrm x\sqrt[3]{3\mathrm x}\right)\cdot\log_3\mathrm x^3}$\par $+\sqrt{\left(\log_3\sqrt[3]{\mathrm x/3}+\log_\mathrm x\sqrt[3]{3/\mathrm x}\right)\cdot\log_3\mathrm x^3}=2$}\par
is equal to
\begin{qoptions}
\opt{A}{4}
\opt{B}{5}
\opt{C}{6}
\opt{D}{7}
\end{qoptions}
\qn{40} Let f(x) be a polynomial function of 9 degree and having leading coefficient 2009.\par
If f$\left(\dfrac1{\mathrm r}\right)=\dfrac1{\mathrm r}$ for r = 1, 2, 3, \ldots\ldots, 9 and f$\left(\dfrac1{10}\right)=\dfrac{\mathrm m}{\mathrm n}$, then $\dfrac{\mathrm n-10\mathrm m}{10}$ is equal to\par
(m , n are coprime)
\begin{qoptions}
\opt{A}{0}
\opt{B}{2010}
\opt{C}{20090}
\opt{D}{none of these}
\end{qoptions}
\qn{41} Let A denotes the sum of the roots of the equation $\dfrac1{5-4\log_4\mathrm x}+\dfrac4{1+\log_4\mathrm x}=3$.\par
B denotes the value of the product of m and n, if 2$^\mathrm m$ = 3 and 3$^\mathrm n$ = 4.\par
C denotes the sum of the integral roots of the equation $\log_{3\mathrm x}\left(\dfrac3{\mathrm x}\right)+\left(\log_3\mathrm x\right)^2=1$.\par
Which of the following alternative(s) is/are correct ?
\begin{qoptions}
\opt{A}{The value of A + B equals 8}
\opt{B}{The value of B + C equals 6}
\opt{C}{The value of A + C $\div$ B equals 8}
\opt{D}{The value of A + C $\div$ B equals 5}
\end{qoptions}
\qn{42} If the complete set of x satisfying $\log_{\log_2\left(\frac{\mathrm x}2\right)}\left(\mathrm x^2-10\mathrm x+22\right)>0$ is $(\mathrm a,\mathrm b)\cup(\mathrm c,\infty)$, then which of the following alternative(s) is/are incorrect ?
\begin{qoptions}
\opt{A}{a + b = 8 $-$ $\sqrt3$}
\opt{B}{a + c = 8 + $\sqrt3$}
\opt{C}{b + c = 10}
\opt{D}{a + c = 10}
\end{qoptions}
\qn{43} Let f(x) = $-5\mathrm x^3+\mathrm{ax}^2+\mathrm{bx}+\mathrm c$ where f(1) = 1, f(2) = 8, f(3) = 27, f(4) = $\lambda$. Also $\mu=\dfrac{\log_224}{\log_{96}2}-\dfrac{\log_2192}{\log_{12}2}$ then which of the following alternative(s) is/are incorrect ?
\begin{qoptions}
\opt{A}{$\lambda-\mu$ = 25}
\opt{B}{$\lambda+\mu$ = 37}
\opt{C}{$\lambda>30$}
\opt{D}{$\mu>3$}
\end{qoptions}
\qn{44} If $\left(\dfrac{\sqrt{\mathrm a+4}}{1-\mathrm a}-1\right)\mathrm x^4-5<0$ for all $\mathrm x\in\mathrm R$, then the complete set of values of a is $[\mathrm m,\mathrm n]\cup(\mathrm p,\infty)$.\par
Which of the following alternative(s) is/are correct ?
\begin{qoptions}
\opt{A}{4m + (3 $-$ 2n)$^2$ $-$ 3p = 2}
\opt{B}{2m + (3 $-$ 2n)$^2$ $-$ p = 13}
\opt{C}{(3 $-$ 2n)$^2$ + m = 17}
\opt{D}{(3 $-$ 2n)$^2$ + p = 21}
\end{qoptions}
\qn{45} If the inequality, $2\log_{0.5}y^2-3+2\mathrm x\log_{0.5}y^2-\mathrm x^2>0$ is valid for atleast one real value of `x', then $y$ can have the value
\begin{qoptions}
\opt{A}{$\cot\left(\dfrac{29\pi}{15}\right)$}
\opt{B}{$\tan\left(\dfrac{23\pi}8\right)$}
\opt{C}{$\tan\left(\dfrac{46\pi}{15}\right)$}
\opt{D}{$\cot\left(\dfrac{9\pi}8\right)$}
\end{qoptions}
\qn{46} If 2$^\mathrm x$ = 49$^{\mathrm x-1}$, then x is equal to :
\begin{qoptions}
\opt{A}{$\dfrac{2\log_27}{\log_27-1}$}
\opt{B}{$\dfrac1{1-\log_7\sqrt2}$}
\opt{C}{$\dfrac2{2+\log_7\frac12}$}
\opt{D}{$\dfrac1{1+\log_{49}2}$}
\end{qoptions}
\qn{47} Let a = $\Bigl(\sqrt{\log_2\sqrt[4]{10}+\log_5\sqrt[4]{10}}-$ $\sqrt{\log_2\sqrt[4]{\tfrac52}+\log_5\sqrt[4]{\tfrac25}}\Bigr)$ $\left(\sqrt{\log_25}\right)$,\par
b = $\Bigl(\sqrt{\log_3\sqrt[4]6+\log_2\sqrt[4]6}-$ $\sqrt{\log_3\sqrt[4]{\tfrac23}+\log_2\sqrt[4]{\tfrac32}}\Bigr)$ $\left(\sqrt{\log_32}\right)$.\par
Find antilog of (a + b) to the base 3
\qn{48} Let s be the sum of all the integers which are less than 10 and are satisfying the equation $\left|\mathrm x^2-5\mathrm x+6\right|-9=\left|\mathrm x^2-6\mathrm x\right|+|3-\mathrm x|$, and p be the sum of all the integers which are less than 7 and are satisfying the equation $\left|\mathrm x^2-5\mathrm x+6\right|=\left|\mathrm x^2-\mathrm x\right|-|6-4\mathrm x|$. Find the value of (s + p)
\qn{49} Let a = $\left(\dfrac{\log_3\left(\log_2{}^25\right)}{\left(2\log_9\left(\left(3^{\log_{15}30}\right)\left(5^{\log_{15}2}\right)\right)\right)-1}\right)$, then $4^{\left(\sqrt2\right)^\mathrm a}$ is equal to
\qn{50} Find the absolute value of sum of all the integral values of k for which the equation $|\mathrm x+1|-|\mathrm x|+3|\mathrm x-1|-2|\mathrm x-2|-2=\mathrm x+\mathrm k$ has more than two solutions.
\qn{51} If any solution of the inequality, $\dfrac{\log_3\left(\mathrm x^2-3\mathrm x+7\right)}{\log_3(3\mathrm x+2)}<1$ is also a solution of the inequality $\mathrm x^2+(5-2\mathrm a)\mathrm x\leq10\mathrm a$, then the least value of a is $\lambda$. Find the value of (4$\lambda$ $-$ 3)
\qn{52} Let P denotes number of digits in 8$^{12}$ 5$^{35}$ and Q denotes number of zeroes after decimal before a significant figure start in $\left(\dfrac8{27}\right)^{20}$ where log$_{10}$2 = 0.3010 \& log$_{10}$3 = 0.4771. The value of (P + Q) equals.
\qn{53} Let a = $\sqrt{\log_23\cdot\log_212\cdot\log_248\cdot\log_2192+16}-\left(\log_212\cdot\log_248\right)+\log_264$, then find the value of $\displaystyle\sum_{\mathrm k=0}^{\infty}\dfrac{\mathrm a^\mathrm k}{\mathrm a^{\mathrm a^\mathrm k}+1}$
\qn{54} Let a = $\dfrac{\left(\log_23-\log_57\right)\left(\log_23-\log_79\right)}{\left(\log_35-\log_57\right)\left(\log_35-\log_79\right)}$, b = $\dfrac{\left(\log_23-\log_35\right)\left(\log_23-\log_79\right)}{\left(\log_57-\log_35\right)\left(\log_57-\log_79\right)}$\par
and c = $\dfrac{\left(\log_23-\log_35\right)\left(\log_23-\log_57\right)}{\left(\log_79-\log_35\right)\left(\log_79-\log_57\right)}$. Find the value of [a + b + c]\par
( [.] denotes greatest integer function)
\end{multicols}
\end{document}
